\documentclass[10pt]{article}
\input{../../notes-style.tex}

\usepackage{multicol}
\usepackage{array}

\title{Ch.~13 Notes\\\large Additional Algorithms and List Forms}
\author{}
\date{}

\begin{document}
\maketitle
\vspace{-2em}

\begin{multicols}{2}

\subsection*{Bubble sort}
Repeatedly pass through the array swapping adjacent
out-of-order pairs.
\begin{lstlisting}[language=C++]
for (int i = 0; i < n - 1; ++i) {
    bool swapped = false;
    for (int j = 0; j + 1 < n - i; ++j)
        if (a[j] > a[j+1]) {
            std::swap(a[j], a[j+1]);
            swapped = true;
        }
    if (!swapped) break;   // early exit if sorted
}
\end{lstlisting}
$O(n^2)$ worst, $O(n)$ best (sorted). Stable, in place.
\textbf{Never use in production} -- insertion sort beats it in
every dimension. Useful only as a teaching example and as a
simple way to ``almost-sort'' partially-sorted data.

\subsection*{Quickselect (find the $k$-th smallest)}
Quicksort, but only recurse into the partition that contains $k$.
\begin{lstlisting}[language=C++]
int quickselect(vector<int>& a, int lo, int hi, int k) {
    int p = lomutoPartition(a, lo, hi);
    if (p == k) return a[p];
    if (k < p) return quickselect(a, lo,    p-1, k);
    else       return quickselect(a, p+1, hi,    k);
}
\end{lstlisting}
Expected $O(n)$; worst $O(n^2)$ (bad pivot). C++:
\texttt{std::nth\_element(first, nth, last)} -- $O(n)$ expected,
puts the $n$-th smallest at the \texttt{nth} position, with
everything smaller before and larger after.

\subsection*{Applications of quickselect}
Median, top-$k$, percentile queries without a full sort. For
streaming / huge $n$: combine with a heap of size $k$.

\subsection*{Bucket sort}
Assumes a roughly uniform distribution over a bounded range.
\begin{enumerate}
    \item Split the range into $k$ buckets.
    \item Distribute items into buckets by a key function.
    \item Sort each bucket (insertion / \texttt{std::sort}).
    \item Concatenate.
\end{enumerate}
Average $O(n + k)$ if items distribute uniformly;
worst $O(n^2)$ (all into one bucket).

\subsection*{Bucket / counting / radix family}
\begin{tabular}{@{}l|l@{}}
\textbf{use when} & \textbf{algo}\\
\hline
small integer range & counting sort\\
uniform over $[0,1)$ or similar & bucket sort\\
multi-digit / multi-byte keys  & radix (LSD / MSD)\\
general / comparisons only     & merge / quick / heap
\end{tabular}

\columnbreak

\subsection*{List ADT vs.\ linked list}
\textbf{List} = sequential positional collection ADT. One
\emph{implementation} is a linked list. Wrap the raw node-pointer
code in a class with invariants (head, tail, size, no cycles).

\begin{lstlisting}[language=C++]
template <class T>
class List {
    struct Node { T val; Node* next; };
    Node* head_ = nullptr;
    Node* tail_ = nullptr;
    std::size_t size_ = 0;
public:
    void push_front(T); void push_back(T);
    void pop_front();   void pop_back();
    std::size_t size() const { return size_; }
    // iterators, destructor, rule of 5 ...
};
\end{lstlisting}

\subsection*{Why wrap, why not just \texttt{std::list}}
\textbf{Wrap own} when learning, when teaching, or when you need
a special invariant (XOR list, intrusive list, lock-free).
\textbf{Prefer \texttt{std::list}} in real code -- it's been tuned,
supports iterators, has the rule of 5 right. \textbf{But usually}:
prefer \texttt{std::vector} unless profiling proves otherwise.
Linked lists lose on cache locality, most of the time.

\subsection*{Circular lists}
Tail's \texttt{next} points to head (singly); plus head's
\texttt{prev} points to tail (doubly).
\begin{itemize}
    \item No null end -- termination is \textbf{``came back to the start''}.
    \item Natural for round-robin schedulers, ring buffers,
    infinite playlists, Josephus problem.
    \item Traversal trap: forgetting the terminator and looping
    forever.
\end{itemize}

\subsection*{Sentinel DLL (the clean variant)}
Allocate a dummy \texttt{head} and \texttt{tail} (or a single
sentinel that points to itself). Real nodes sit between.
\begin{itemize}
    \item No empty-list special case.
    \item Insert / remove: always 4 pointer writes, never a nullptr check.
    \item Matches \texttt{std::list}'s internal shape.
\end{itemize}

\subsection*{When \emph{not} to reach for a linked list}
\begin{itemize}
    \item Random access / indexing by position (vector wins).
    \item Dense numeric data with tight loops (cache murder).
    \item Tiny elements where the pointer overhead dominates.
    \item When you need \texttt{sort} / \texttt{binary\_search} --
    no random access, no \texttt{std::sort}.
\end{itemize}

\subsection*{Top gotchas}
\begin{itemize}
    \item Using bubble sort in anything that ships.
    \item Quickselect with first-element pivot on sorted data --
    quadratic; shuffle or use random pivot.
    \item Bucket sort with skewed data -- degenerates to one
    bucket's sort.
    \item Forgetting the rule of 5 in a hand-rolled \texttt{List}
    (copy / move / destructor leaks).
    \item Circular-list traversal with \texttt{while (n)} instead
    of \texttt{do \{...\} while (n != start)}.
\end{itemize}

\end{multicols}
\end{document}
